%Dokumentenklasse
\documentclass[
	a4paper,										%Format
	12pt,											%Schriftgröße
	twoside,										%Zweiseitig
	headsepline,								%horizontale Linie oben
	footsepline,									%horizontale Linie unten
	captions=tableheading,					%Überschrift über Tabelle (nicht unter)
]{scrartcl}
\brokenpenalty=1000000
%Pakete einbinden
\usepackage{floatrow}
\floatsetup[table]{style=plaintop}
\floatsetup[longtable]{LTcapwidth=table}
\usepackage{overcite}
\usepackage{siunitx}
\usepackage[ngerman]{babel} 			%Neue deutsche Rechtschreibung
\usepackage{grffile}
\usepackage{graphicx} 					%Bilder einbinden
\usepackage{xcolor}						%Farben für Schrift und Plots
\usepackage{pgfplots}					%Plots und Tabellen aus Datensätzen.
\usepackage{amsmath}						%Mathe
\usepackage{amssymb}					%Mathesymbole
\usepackage{longtable}					%Seitenübergreifende Tabellen zulassen
\usepackage{tabularx}
\usepackage{geometry}					%Seitenränder
\usepackage{ae}								%skalierbare Schriften (im PDF-File)
\usepackage[activate]{pdfcprot}		%optischer Randausgleich
\usepackage{setspace}					%Zeilenabstand: \onehalfspacing
\usepackage[utf8]{inputenc}				%Zeichenkodierung für Umlaute
\usepackage{scrlayer-scrpage}
\usepackage{mhchem}						%Chemische Formeln (mit \ce )
\usepackage{chemfig}
\usepackage{caption}										%Überschriften formatieren
\captionsetup{labelfont={bf,it},tablename=Tab.,figurename=Abb.}				%Überschriften formatieren
\captionsetup[table]{singlelinecheck=false,justification=raggedright}
\captionsetup[figure]{justification=raggedright}				

%Abbildungs und Tabellennummerierung anpassen
\renewcommand*\figureformat {\figurename~\arabic{section}.\arabic{subsection}.\thefigure.}	%Punkt nach Abb. #
\renewcommand*\tableformat {\tablename~\arabic{section}.\arabic{subsection}.\thetable:}		%Doppelpunkt nach Tab. #
\renewcommand*\captionformat {~}								%Nur Abstand, kein Zeichen nach Tab. #: bzw Abb. #. 

%Literatur als Subsubsection definieren und von Literaturverzeichnis in Literatur umbenennen
\makeatletter
\renewcommand*\bib@heading{\section{\bibname}			
\@mkboth{\sectionmarkformat \bibname}{\sectionmarkformat \bibname}}
\makeatother
\addto\captionsngerman{\renewcommand*\bibname{Literatur}} 


\usepackage{prettyref}						%Nummerierung der Tabellen und Abbildungen im Text an Tab. x.x.# anpassen
%%% Für Abbildungen %%%
\newrefformat{fig}{Abb.~\arabic{section}.\arabic{subsection}.\ref{#1}.}
%%% Für Tabellen %%%
\newrefformat{tab}{Tab.~\arabic{section}.\arabic{subsection}.\ref{#1}}


%Links im Dokument:
\usepackage[									%als letztes Paktet einbinden!
	linkcolor=black,							%Farbe der Links
	citecolor=black,							%Farbe der Link-Ziele
	urlcolor=black,								%Farbe von URLs
	colorlinks=true,							%Links nicht als Kasten, sondern direkt als Text anzeigen
]{hyperref}										%Links im PDF-File, unterstützung von URLs mit /url{}
\renewcommand\citeform[1]{[#1]}
%Dokumentformatierungen
\usepackage{titlesec}

\titleformat{\section}
	{\normalfont\Large\bfseries}{\arabic{section}.}{1.5em}{}
\titleformat{\subsection}
	{\normalfont\Large\bfseries}{\arabic{section}.\arabic{subsection}.}{1.0em}{}
\titleformat{\subsubsection}
	{\normalfont\Large\bfseries}{\arabic{section}.\arabic{subsection}.\arabic{subsubsection}.}{0.75em}{}

\setkomafont{sectioning}{\bfseries} 	%Serifenschrift in Überschriften
\setkomafont{subsubsection}{\large}
\pagenumbering{arabic}					%arabische Seitenzahlen
\geometry{left=3cm, right=2cm, top=2.5cm, bottom=2cm} 	%Seitenränder
\parindent 0pt 									%Absätze einrücken um 0Pt
\parskip 10pt									%Abstand nach Absatz
\onehalfspacing								%Zeilenabstand
\pagestyle{scrheadings}					%Kopf- und Fussleistenformatierung
\headsep0.4cm								%Abstand Kopfleiste-Text
\footskip0.8cm									%Abstand Fussleiste-Text
\urlstyle{same}								%URLs normal formatieren, nicht als Typewriter
\raggedbottom

%Tabellenkram


%Befehle definieren

%%%%% create chemequation environment

\newcounter{chemequation}
\renewcommand*\thechemequation{(\Roman{chemequation})}
\newenvironment{chemequation}{%
	\stepcounter{chemequation}%
	\begin{equation}}%
{\tag*{\thechemequation}%
	\end{equation}}

%%%%% create chemtag to iterate and ref chemequations in an align environment

\newcommand*\chemtag{%
	\stepcounter{chemequation}%
	\tag*{\thechemequation}}
	
	\counterwithin*{equation}{section}
\counterwithin*{chemequation}{section}
\counterwithin*{table}{subsection}
\counterwithin*{figure}{subsection}

%%%%% to create a chemicalequation float

\usepackage{float}
\newfloat{chemicalequation}{htbp}{cex}%[chapter]
\floatname{chemicalequation}{}

\newcommand{\zitat}[1]{\textsuperscript{[\cite{#1}]}}

\newcommand{\bild}[4]{					%Bilder mit Beschriftung einfügen
	\begin{figure}[#4]						%#4: Platzierung (htbp)
		\centering								%Bilder mittig einfügen
		\captionbox								%Beschriftungsumgebung
			{#2\label{fig: #1}}				%#2: Beschriftung
			{\includegraphics[#3]{#1}}	%#1: Dateiname, #3: Bildeigenschaften
	\end{figure}}

\newcommand{\formel}[2]{				%Formel einfügen und Referenz erzeugen
	\begin{equation}
		\label{eq: #1}							%#1: Formelname
		#2											%Formelinhalt
	\end{equation}}

\newcommand{\cformel}[2]{				%Chemische Formel einfügen und Referenz erzeugen
	\begin{chemequation}
		\label{ceq: #1}					%#1: Formelname (optional)
		\ce{#2}							%#2: Formelinhalt
	\end{chemequation}\\}

\newcommand{\ox}[2]{
	\overset{\text{#2}}{\ce{#1}}}
	
\newcommand{\tabelle}[6]{				%Tabelle mit Beschriftung einfügen
\begin{table}[#3]							%#3: Platzierung (htbp)
\centering										%Tabelle mittig einfügen
\captionbox										%Beschriftungsumgebung
{#2												%#2: Tabellenbeschriftung
\label{tab: #1}}								%#1: Tabellenname
{\begin{tabular}{#4}						%#4: Anzahl und Ausrichtung Spalten (l c p{xcm} r)     Bei p-Spalten wird Text atomatisch umgebrochen
\hline 											%horizontale Linie
#5													%#5: Obere Zeile
\hline
#6													%#6: Datenzeilen
\hline 
\end{tabular}}
\end{table}}

\newcommand{\langetabelle}[6]{  %Seitenübergreifende Tabelle, Kopfzeile wird 
	\begin{centering}           %auf jeder neuen Seite ebenfalls angezeigt
		{\begin{longtable}[#3]{#4}
				\caption{#2\label{tab: #1}}\\				
				\hline
				#5
				\hline
				\endfirsthead
				\hline
				#5
				\hline
				\endhead
				#6												
				\hline 
		\end{longtable}}
\end{centering}}

%Plots erstellen&&&&&&&&

%#1: Referenz
%#2: Name
%#3: x label
%#4: y label
%#5: x-werte,y-werte (Spaltenweise)
%#6: Quelle (.txt Datei)
%#7: "Nummerierung" (manuell) <.<

\newcommand{\plot}[6]{
	\begin{figure}[H]	
		\centering
		\begin{tikzpicture}
	
		\begin{axis}[
		title style={at={(0.5,0)},
		anchor=north,
		yshift=-40},
		xlabel={#3},
		ylabel={#4},
		legend pos=south east,
		legend entries={},
		]
		\addplot[color=black,mark=*,only marks] table [#5] {#6};
		\label{plot: plot1}
		\end{axis}
		\end{tikzpicture}
		\caption{#2}
		\label{fig: #1}
	\end{figure}
}

\newcommand{\Rn}[1]{\romannumeral#1\relax}                   %\Rn{1} = i
\newcommand{\RN}[1]{\uppercase\expandafter{\romannumeral#1}} %\RN{1} = I

%Dokument:%%%%%%%%%%%%%%%%%%%%%%%%%%%%%%%%%%%%%%%%%%%%%%%%%%%%%%%%%%%%%%%%%%%%%%%%%%%%%
\begin{document}
\thispagestyle{empty}
	
%Titelseite
\begin{titlepage}
	\begin{center}
		\hrule													%waagerechte Linie		\vspace*{5cm}		\Huge {\bf {Titel}}\\							%Titel		\vspace*{0.5cm}		\large {\bf {Untertitel}}\\						%Untertitel
		\vspace*{\fill}
		\Huge {\textbf {Rasterelektronenmikroskopie und EDX-Elementanalyse}}\\
		\vspace*{1cm}
		\normalsize {\textbf {Praktikum Nanostrukturphysik}}\\
		\vspace*{1cm}
		\large {\textbf{Victor Kemme, Kilian Krötzsch}}\\
		\normalsize{Gruppe 9}								%Gruppe
		\vfill
	\end{center}
	\begin{tabbing}
		\hspace{3cm}\=\kill
		Assistenten: \> Sarah Scheitz, Dr. Robert Frömter \\						%Assistent
		Version: \> 1\\										%Version
		Datum: \> \today \\								%Datum
	\end{tabbing}
	\hrule														%Waagerechte Linie
\end{titlepage}
\setcounter{section}{0}								%Versuchstag #.x.x
\setcounter{subsection}{0}							%Versuchsnummer x.#.x
\normalsize
\setcounter{page}{1}

\section{Einleitung und Aufgabenstellung}

\section{Theoretischer Hintergrund}

\section{Experimenteller Teil}

\section{Ergebnisse und Diskussion}

\section{Zusammenfassung}

\section{Sicherheit und Entsorgung}

\begin{thebibliography}{}
	\bibitem{MC}
	PD Dr. Ulrich Handge; Prof. Dr. Gerrit Albert Luinstra; Dr. Axel Thomas Neffe und Dr. Werner Pauer im \textit{Skript zur Vorlesung 2 im Modul makromolekulare Chemie}, Universität Hamburg, \textbf{2019}
	\bibitem{Praktikumsskript S.34}
	H.Heller im \textit{Praktikumsskript zum Praktikum „Grundlagenpraktikum der Chemie für
	Studierende der Nanowissenschaften“}, Universität Hamburg, Hamburg, \textbf{2019}, S. 34.
	\bibitem{xxx}
	H.Heller im \textit{Seminarfoliensatz 5 zum Praktikum „Grundlagenpraktikum der Chemie für
	Studierende der Nanowissenschaften“}, Universität Hamburg, Hamburg, \textbf{2019}, S. 8ff.
	
\end{thebibliography}

\formel{mom}{g^{(1)}(\tau) = \sqrt{g^{(2)}(\tau) - 1}}
\formel{}{\ln(g^{(1)}) = -\Gamma \tau = -q^2 D_0 \tau}
\formel{}{R_H = \frac{-k T 8 \pi n^2 \sin^2({\frac{\theta}{2}})}{3 \eta m \lambda^2}}
\formel{}{D_0 = \frac{- m \ \lambda^2}{16 \ \pi^2 \ n^2 \ \sin^2(\frac{\theta}{2})}}

aus \ref{eq: mom} folgt
aus \prettyref{eq: mom} folgt

\formel{}{\phi = k(b-a) = \vec{r} \cdot \vec{k}_S - \vec{r} \cdot \vec{k}_0}
\formel{}{|\vec{k}_0| = |\vec{k}_S| = k = \frac{2\pi}{\lambda}}
\formel{}{|\vec{q}| = q = \frac{4\pi}{\lambda}\sin{\frac{\theta}{2}},  \ \text{mit} \ \lambda = \frac{\lambda_0}{n}}


\formel{}{g^{(2)}(\tau) \lim \limits_{\text{T} \to \infty} \frac{1}{\text{T}} \int_{0}^{\text{T}} I(t) \cdot I(t+\tau) dt = \langle I(t) \cdot I(t+\tau) \rangle}
\formel{}{\lim \limits_{\tau \to 0} \langle I(t) \cdot I(t+\tau) \rangle = \langle I^2(t) \rangle}
\formel{}{\lim \limits_{\tau \to \infty} \langle I(t) \cdot I(t+\tau) \rangle = \langle I(t) \rangle^2}

\formel{}{v_m = \sqrt{\frac{2 \cdot k \cdot T}{m}}}

\formel{}{I(t) = \vec{E}(t) V^2}
\formel{}{g^{(1)}(\tau) = \langle E_S(t) \cdot E_S(t+ \tau) \rangle}
\formel{}{g^{(2)}(\tau) = 1 + \bigg(g^{(1)}(\tau) \bigg)^2}
\formel{}{g_n^{(1)}(\tau) = \frac{ \langle E_S(t) \cdot E_S(t+ \tau) \rangle}{\langle E_S(t) \cdot E_S(t) \rangle} = \frac{ \langle E_S(t) \cdot E_S(t+ \tau) \rangle}{\langle I(t) \rangle} }
\formel{}{g_n^{(1)}(\tau) = e^{-\Gamma \tau} = e^{-q^2D_0 \tau}}
\formel{}{D_0= \frac{kT}{6 \pi \eta R_H}}
\formel{}{D = \frac{kT}{f}}
\formel{}{f = 6 \pi \eta R_H}

\formel{}{g^{(1)}(\tau) = \int_{0}^{\infty} G(\Gamma) \cdot e^{-\Gamma \tau} d\Gamma , \ \text{mit} \int_{0}^{\infty} G(\Gamma) d\Gamma = 1}
\formel{}{\ln{\bigg(g^{(1)}(\tau)\bigg)} = \sum_{n = 1}^{\infty} \Gamma_n \frac{(-\tau)^n}{n!} = -\Gamma_1 \tau + \frac{\Gamma_2 \tau^2}{2} \pm \mathcal{O}(\tau^n)}
\formel{}{\Gamma_n = \frac{\partial^n \ln{\big(g^{(1)}(\tau)\big)}}{\partial (-\tau)^n} \Bigg|_{\tau = 0}}
\formel{}{\sigma^2 = \frac{\langle (\Gamma - \langle \Gamma \rangle )^2 \rangle}{\langle \Gamma \rangle}}

\end{document}


% <- Für Kommentare oder Code der gerade nicht gebraucht wird, wird beim kompilieren nicht beachtet
%Um Bilder einzufügen: \bild{'Dateiname' (Datei muss im gleichem Ordner mit der TeX Datei liegen)}{'Bildunterschrift'}{width= ...cm (Größe des Bildes)}{H (Ort des Bildes, H=Hier)}

%Für Summenformeln im Text: \ce{z.B.: H2SO4 (Zahlen werden automtisch tief gestellt)}

% Um Formeln im Text zu referieren: \prettyref{fig (für Bilder und Graphen)/ eq (für mathematische Formeln)/ ceq (für chemische Formeln): 'Referenz'}

%Für chemische Formeln: \cformel['Referenz']{2\ox{Fe}{+\RN{3}}Cl3 + \ox{Fe}{+\RN{2}}Cl2 + 8NH4+ + 8 OH- \ce{<=>[{N2H4}]} \ox{Fe}{+\RN{2}}O\cdot \ox{Fe}{+\RN{3}}2O3 \ce{v} + 4 H2O + 8NH4+ + 4 Cl-}

%Für mathematische Formeln: \formel{'Referenz'}{2\cdot d\cdot \ce{sin}(\theta)= n \lambda}

%Tabellen: Online Converter um Exel Tabellen in LaTeX Format zu bringen: https://www.tablesgenerator.com
%\tabelle{'Referenz'}{'Tabellenüberschrift'}{H}{'Breite der Spalte bzw. Textausrichtung' (Spaltenbreite passt sich Textlänge an: l=links, r=rechts, c=zentriert. Spaltenbreite fest und Blocksatz: p{...cm}), Jede Spalte muss einzeln definiert werden, z.B l l l = 3 Spalten}{'Spaltenüberschrift' muss der Anzahl an Spalten entsprechen. Spaltenüberschriften müssen durch ein & getrennt werden. Außerdem wichtig am Ende: \\}{'Tabelleneinträge', müssen ebenfalls durch ein & getrennt werden und am Ende der Zeile ein \\ haben
%	- & - & - & \\
%	- & - & - & \\}
%Beispiel für eine Tabelle (zum Ausführen '%' entfernen):
%\tabelle{Auflösungsvermögen}{Experimentell ermitteltes und theoretisches Auflösungsvermögen für Objektiv 20X und 63X bei schräger und gerader Beleuchtung.}{H}{l l p{2cm} p{3cm}  p{3cm}}{Objektiv&Beleuchtung& Spalt mit einem Peak&experimentelles Auflösungsvermögen~/~nm&theoretisches Auflösungsvermögen~/~nm\\}
%	{20X&gerade&16&>750&1125\\
%	20X&schräg&9&690&563\\
%	63X&gerade&11&>750&600\\
%	63X&schräg&7&584&300\\}

